%USEPACKAGES

%package list
\documentclass{article}
\usepackage[top=3cm, bottom=3cm, outer=3cm, inner=3cm]{geometry}
\usepackage{multicol}
\usepackage{graphicx}
\usepackage{url}
%\usepackage{cite}
\usepackage{hyperref}
\usepackage{array}
%\usepackage{multicol}
\newcolumntype{x}[1]{>{\centering\arraybackslash\hspace{0pt}}p{#1}}
\usepackage{natbib}
\usepackage{pdfpages}
\usepackage{multirow}
\usepackage[normalem]{ulem}
\useunder{\uline}{\ul}{}
\usepackage{svg}
\usepackage{xcolor}
\usepackage{listings}
\lstdefinestyle{ascii-tree}{
    literate={├}{|}1 {─}{--}1 {└}{+}1 
  }
\lstset{basicstyle=\ttfamily,
  showstringspaces=false,
  commentstyle=\color{red},
  keywordstyle=\color{blue}
}
%\usepackage{booktabs}
\usepackage{caption}
\usepackage{subcaption}
\usepackage{float}
\usepackage{array}
\usepackage[utf8]{inputenc}
\newcolumntype{M}[1]{>{\centering\arraybackslash}m{#1}}
\newcolumntype{N}{@{}m{0pt}@{}}


%Lab data
\newcommand{\itemEmail}{jariasq@unsa.edu.pe}
\newcommand{\itemStudent}{Jhonatan David Arias Quispe}
\newcommand{\itemCourse}{Fundamentos de Programacion 2}
\newcommand{\itemCourseCode}{1701213}
\newcommand{\itemSemester}{II}
\newcommand{\itemUniversity}{Universidad Nacional de San Agustín de Arequipa}
\newcommand{\itemFaculty}{Facultad de Ingeniería de Producción y Servicios}
\newcommand{\itemDepartment}{Departamento Académico de Ingeniería de Sistemas e Informática}
\newcommand{\itemSchool}{Escuela Profesional de Ingeniería de Sistemas}
\newcommand{\itemAcademic}{2023 - B}
  



\newcommand{\itemInput}{Del 4 Enero 2024}
\newcommand{\itemOutput}{Al 11 Enero 2024}
\newcommand{\itemPracticeNumber}{20}
\newcommand{\itemTheme}{Laboratorio 20}
\newcommand{\itemUrl}{\url{https://github.com/JhonatanDczel/fp2-23b/tree/main/fase03/lab\itemPracticeNumber}}


%Format Document
\usepackage[english,spanish]{babel}
\usepackage[utf8]{inputenc}
\AtBeginDocument{\selectlanguage{spanish}}
\renewcommand{\figurename}{Figura}
\renewcommand{\refname}{Referencias}
\renewcommand{\tablename}{Tabla} %esto no funciona cuando se usa babel
\AtBeginDocument{%
	\renewcommand\tablename{Tabla}
}

\usepackage{fancyhdr}
\pagestyle{fancy}
\fancyhf{}
\setlength{\headheight}{30pt}
\renewcommand{\headrulewidth}{1pt}
\renewcommand{\footrulewidth}{1pt}
\fancyhead[L]{\raisebox{-0.2\height}{\includegraphics[width=3cm]{img/logo_episunsa.png}}}
\fancyhead[C]{\fontsize{7}{7}\selectfont	\itemUniversity \\ \itemFaculty \\ \itemDepartment \\ \itemSchool \\ \textbf{\itemCourse}}
\fancyhead[R]{\raisebox{-0.2\height}{\includegraphics[width=1.2cm]{img/logo_abet}}}
\fancyfoot[L]{Jhonatan David Arias Quispe}
\fancyfoot[C]{\itemCourse}
\fancyfoot[R]{Página \thepage}


% Code Format

\usepackage{listings}
\usepackage{color, colortbl}
\definecolor{dkgreen}{rgb}{0,0.6,0}
\definecolor{gray}{rgb}{0.5,0.5,0.5}
\definecolor{mauve}{rgb}{0.58,0,0.82}
\definecolor{codebackground}{rgb}{0.95, 0.95, 0.92}
\definecolor{tablebackground}{rgb}{0.8, 0, 0}

\lstset{frame=tb,
	language=bash,
	aboveskip=3mm,
	belowskip=3mm,
	showstringspaces=false,
	columns=flexible,
	basicstyle={\small\ttfamily},
	numbers=none,
	numberstyle=\tiny\color{gray},
	keywordstyle=\color{blue},
	commentstyle=\color{dkgreen},
	stringstyle=\color{mauve},
	breaklines=true,
	breakatwhitespace=true,
	tabsize=3,
	backgroundcolor= \color{codebackground},
}
\usepackage{color}
    \definecolor{bat-bg}{rgb}{.16,.19,.2}
    \definecolor{bat-norm}{rgb}{.9,.9,.9}
    \definecolor{bat-comm}{rgb}{.4,.45,.48}
    \definecolor{bat-kw}{rgb}{.57,.78,.39}
    \definecolor{bat-str}{rgb}{.2,.38,.62}
    \definecolor{arch-celeste}{rgb}{0.263, 0.765, 0.961}
    \definecolor{arch-lime}{rgb}{0.588, 0.773, 0.506}



\usepackage{listings}

\lstset{
    language=java,
    basicstyle=\color{bat-norm}\footnotesize\ttfamily,
    numbers=left,
    numberstyle=\normalsize,
    numbersep=7pt,
    backgroundcolor=\color{bat-bg},  
    commentstyle=\color{bat-comm},
    keywordstyle=\color{arch-celeste},
    numberstyle=\tiny\color{bat-bg},
    stringstyle=\color{arch-lime},
    tabsize=2
}


\begin{document}
%Data table

	
	\vspace*{10pt}
	
	\begin{center}	
		\fontsize{17}{17} \textbf{ Informe de Laboratorio \itemPracticeNumber}
	\end{center}
	\centerline{\textbf{\Large Tema: \itemTheme}}
	\vspace*{0.5cm}	

	\begin{flushright}
		\begin{tabular}{|M{2.5cm}|N|}
			\hline 
			\rowcolor{tablebackground}
			\color{white} \textbf{Nota}  \\
			\hline 
			     \\[30pt]
			\hline 			
		\end{tabular}
	\end{flushright}	

	\begin{table}[H]
		\begin{tabular}{|x{4.9cm}|x{4.3cm}|x{5.1cm}|}
			\hline 
			\rowcolor{tablebackground}
			\color{white} \textbf{Estudiante} & \color{white}\textbf{Escuela}  & \color{white}\textbf{Asignatura}   \\
			\hline 
			{\itemStudent \par \itemEmail} & \itemSchool & {\itemCourse \par Semestre: \itemSemester \par Código: \itemCourseCode}     \\
			\hline 			
		\end{tabular}
	\end{table}		
	
	\begin{table}[H]
		\begin{tabular}{|x{4.7cm}|x{4.8cm}|x{4.8cm}|}
			\hline 
			\rowcolor{tablebackground}
			\color{white}\textbf{Laboratorio} & \color{white}\textbf{Tema}  & \color{white}\textbf{Duración}   \\
			\hline 
			\itemPracticeNumber & \itemTheme & 04 horas   \\
			\hline 
		\end{tabular}
	\end{table}
	
	\begin{table}[H]
		\begin{tabular}{|x{4.7cm}|x{4.8cm}|x{4.8cm}|}
			\hline 
			\rowcolor{tablebackground}
			\color{white}\textbf{Semestre académico} & \color{white}\textbf{Fecha de inicio}  & \color{white}\textbf{Fecha de entrega}   \\
			\hline 
			\itemAcademic & \itemInput &  \itemOutput  \\
			\hline 
		\end{tabular}
	\end{table}

	\section{Actividades}
	\begin{itemize}		
      \item Cree un Proyecto llamado Laboratorio7
      \item Usted deberá crear las dos clases Soldado.java y VideoJuego4.java. Puede reutilizar lo desarrollado en Laboratorios anteriores.
      \item Del Soldado nos importa el nombre, puntos de vida, fila y columna (posición en el tablero).
      \item El juego se desarrollará en el mismo tablero de los laboratorios anteriores. Para el tablero utilizar la estructura de datos más adecuada.
      \item Tendrá 2 Ejércitos (utilizar la estructura de datos más adecuada). Inicializar el tablero con n soldados aleatorios entre 1 y 10 para cada Ejército. Cada soldado tendrá un nombre autogenerado: Soldado0X1, Soldado1X1, etc., un valor de puntos de vida autogenerado aleatoriamente [1..5], la fila y columna también autogenerados aleatoriamente (no puede haber 2 soldados en el mismo cuadrado). Se debe mostrar el tablero con todos los soldados creados y sus puntos de vida (usar caracteres como y otros y distinguir los de un ejército de los del otro ejército). Además de los datos del Soldado con mayor vida de cada ejército, el promedio de puntos de vida de todos los soldados creados por ejército, los datos de todos los soldados por ejército en el orden que fueron creados y un ranking de poder de todos los soldados creados por ejército (del que tiene más nivel de vida al que tiene menos) usando 2 diferentes algoritmos de ordenamiento. Finalmente, que muestre qué ejército ganará la batalla (indicar la métrica usada para decidir al ganador de la batalla). Hacer el programa iterativo.
	\end{itemize}
		
	\section{Equipos, materiales y temas utilizados}
	\begin{itemize}
		\item Sistema Operativo ArchCraft GNU Linux 64 bits Kernell
		\item NeoVim
		\item OpenJDK 64-Bit 20.0.1 
		\item Git 2.42.0
		\item Cuenta en GitHub con el correo institucional.
		\item Programación Orientada a Objetos.
		\item Creacion de programas con CLI	
            \item Bilioteca Graphics (origen propio)
	\end{itemize}
	\section{URL de Repositorio Github}
	\begin{itemize}
            \item URL del Repositorio GitHub para clonar o recuperar.
            \item \url{https://github.com/JhonatanDczel/fp2-23b.git}
            \item URL para el laboratorio \itemPracticeNumber{} en el Repositorio GitHub.
            \item \itemUrl
	\end{itemize}




\section{Actividad 1}
\subsection{Para las pruebas de ejecucion}
\begin{itemize}
    \item Para la prueba de ejecucion de cada uno de los metodos que desarrollaremos, vamos a usar las siguientes naves:
\end{itemize}
\begin{lstlisting}
    Naves creadas:
    
    Nave 1:
    Nombre: cesar
    Estado: true
    Puntos: 23
    
    Nave 2:
    Nombre: Alvaro
    Estado: true
    Puntos: 6
    
    Nave 3:
    Nombre: Bannia
    Estado: true
    Puntos: 89
\end{lstlisting}

\section{Completando los metodos faltantes}
\begin{itemize}
    \item La acitivdad del laboratorio 04 nos pide crear los siguientes metodos:
\end{itemize}
\subsection{Busqueda Lineal}
\begin{itemize}
    \item La busqueda lineal es una busqueda simple, en la que se recorre todo el arreglo en busca de una informacion dada
\end{itemize}
\begin{lstlisting}[language=java, caption={Metodo de busqueda lineal}]
      public static int busquedaLinealNombre(Nave[] flota, String s){
        for(int i = 0; i < flota.length; i++){
          if(flota[i].getNombre().equals(s))
            return i;
        }
        return -1;
      } 
\end{lstlisting}

\begin{itemize}
    \item El metodo devuelve la posicion del elemento en caso encontrarlo
    \item Y de lo contrario devuelve -1 lo que indica que el elemento no se encontro
\end{itemize}
\subsubsection{Ejecucion}
\begin{itemize}
    \item Ahora veremos la prueba del metodo:
\end{itemize}
    \begin{lstlisting}[language=bash, caption={Prueba del metodo}]
    Busqueda lineal, inserte un nombre:
    Bannia
    Nombre: Bannia
    Estado: true
    Puntos: 89
    
\end{lstlisting}

\subsection{Ordenamiento por nombre: burbuja}
\begin{itemize}
    \item Ahora veremos como implementar el bubble sort para ordenar las naves con respecto a sus nombres
\end{itemize}
    \begin{lstlisting}[language=java, caption={Metodo de ordenamiento burbuja por nombre}]
      public static void ordenarPorNombreBurbuja(Nave[] flota){
        for(int i = 0; i < flota.length - 1; i++){
          for(int j = 0; j < flota.length - 1 - i; j++){
            if(esMayor(flota[j].getNombre(), flota[j + 1].getNombre()))
              intercambiar(flota, j, j + 1);
          }
        }
      }
\end{lstlisting}
\begin{itemize}
    \item Como vemos, estamos utilizando un metodo que compara dos strings y devuelve cual es mayor lexicograficamente:
\end{itemize}
\begin{lstlisting}[language=java, caption={Sub metodo para comparar dos strings}]
  public static boolean esMayor(String s1, String s2){
    return s1.compareToIgnoreCase(s2) > 0;
  }
\end{lstlisting}
\begin{itemize}
    \item Ahora tenemos una manera de saber el orden de precedencia entre strings
    \item Adicionalmente a esto necesitamos otro metodo para intercambiar elementos de un arreglo de naves:
\end{itemize}
\begin{lstlisting}[language=java, caption={Sub metodo para intercambiar dos elementos}]
      public static void intercambiar(Nave[] flota, int i, int j){
        Nave aux = flota[i];
        flota[i] = flota [j];
        flota[j] = aux;
      }
\end{lstlisting}
\subsubsection{Ejecucion}
\begin{itemize}
    \item Ahora veremos la prueba del metodo:
\end{itemize}
    \begin{lstlisting}[language=bash, caption={Prueba del metodo}]
    Ahora usaremos el algoritmo de ordenamiento burbuja, con respecto a Nombres.
    Naves ordenadas:
    
    
    Nave 1:
    Nombre: Alvaro
    Estado: true
    Puntos: 6
    
    Nave 2:
    Nombre: Bannia
    Estado: true
    Puntos: 89
    
    Nave 3:
    Nombre: cesar
    Estado: true
    Puntos: 23
    
\end{lstlisting}
\subsection{Ordenamiento por puntos: burbuja}
\begin{itemize}
    \item Ahora veremos como ordenar los elementos de un arreglo respecto a sus puntos usando el algoritmo burbuja
    \item Hacer esto es muy similar al ordenamiento por nombre, unicamente cambiando getNombre por getPuntos:
\end{itemize}
\begin{lstlisting}[language=java, caption={Metodo para ordenar por puntos}]
      public static void ordenarPorPuntosBurbuja(Nave[] flota){
        for(int i = 0; i < flota.length - i; i++){
          for(int j = 0; j < flota.length - 1 - i; j++){
            if(flota[j].getPuntos() > flota[j + 1].getPuntos())
              intercambiar(flota, j, j + 1);
          }
        }
      }
\end{lstlisting}
\begin{itemize}
    \item Ahora en lugar de comparar dos strings, directamente se comparan los valores numericos de sus puntos
\end{itemize}
\subsubsection{Ejecucion}
\begin{itemize}
    \item Ahora veremos la prueba del metodo:
\end{itemize}
    \begin{lstlisting}[language=bash, caption={Prueba del metodo}]
    Ahora usaremos el algoritmo de ordenamiento burbuja, con respecto a Puntos.
    Naves ordenadas:
    
    
    Nave 1:
    Nombre: Alvaro
    Estado: true
    Puntos: 6
    
    Nave 2:
    Nombre: cesar
    Estado: true
    Puntos: 23
    
    Nave 3:
    Nombre: Bannia
    Estado: true
    Puntos: 89
     
\end{lstlisting}

\subsection{Ordenamiento por nombre: seleccion}
\begin{itemize}
    \item Ahora veremos el otro algoritmo de ordenamiento, el de seleccion
    \item La estructura basica del algoritmo, es ir seleccionando de izquierda a derecha los elementos mas chicos para ordenar el arreglo
    \item Su implementacion para ordenar una serie de naves lexicograficamente es la siguiente:
\end{itemize}
\begin{lstlisting}[language=java, caption={Metodo para ordenar por nombre}]
      public static void ordenarPorNombreSeleccion(Nave[] flota){
        for(int i = 0; i < flota.length - 1; i++){
          int min = i;
          for(int j = i + 1; j < flota.length; j++){
            if(esMayor(flota[min].getNombre(), flota[j].getNombre()))
              min = j;
          }
          intercambiar(flota, i, min);
        }
      }
\end{lstlisting}
\subsubsection{Ejecucion}
\begin{itemize}
    \item Ahora veremos la prueba del metodo:
\end{itemize}
    \begin{lstlisting}[language=bash, caption={Prueba del metodo}]
    Ahora usaremos el algoritmo de ordenamiento seleccion, con respecto a Nombres.
    Naves ordenadas:
    
    
    Nave 1:
    Nombre: Alvaro
    Estado: true
    Puntos: 6
    
    Nave 2:
    Nombre: Bannia
    Estado: true
    Puntos: 89
    
    Nave 3:
    Nombre: cesar
    Estado: true
    Puntos: 23
    
\end{lstlisting}


\subsection{Ordenamiento por puntos: seleccion}
\begin{itemize}
    \item Ahora usaremos este mismo algoritmo para ordenar el array con respecto a los puntos
    \item Es la misma estructura y el mismo sentido, en lugar de unsar getNombre, usamos getPuntos
\end{itemize}
\begin{lstlisting}[language=java, caption={Metodo para ordenar por puntos}]
      public static void ordenarPorPuntosSeleccion(Nave[] flota){
    for(int i = 0; i < flota.length - 1; i++){
      int min = i;
      for(int j = i + 1; j < flota.length; j++){
        if(flota[j].getPuntos() < flota[min].getPuntos())
          min = j;
      }
      intercambiar(flota, i, min);
    }
  }
\end{lstlisting}
\subsubsection{Ejecucion}
\begin{itemize}
    \item Ahora veremos la prueba del metodo:
\end{itemize}
    \begin{lstlisting}[language=bash, caption={Prueba del metodo}]
    Ahora usaremos el algoritmo de ordenamiento seleccion, con respecto a Puntos.
    Naves ordenadas:
    
    
    Nave 1:
    Nombre: Alvaro
    Estado: true
    Puntos: 6
    
    Nave 2:
    Nombre: cesar
    Estado: true
    Puntos: 23
    
    Nave 3:
    Nombre: Bannia
    Estado: true
    Puntos: 89
      
\end{lstlisting}

\subsection{Ordenamiento por nombre: insercion}
\begin{itemize}
    \item Ahora veremos el algoritmo de insercion
    \item La idea cae en como ordenamos nuestras cartas cuando tenemos una baraja, agarramos la carta y la "insertamos" en su lugar
    \item Vamos a hacer esto pero con codigo, empezamos:
\end{itemize}
\begin{lstlisting}[language=java, caption={Metodo para ordenar por nombre}]
      public static void ordenarPorNombreInsercion(Nave[] flota){
    for(int i = 1; i < flota.length; i ++){
      Nave actual = flota[i];
      int j = i - 1;
      while(j >= 0 && esMayor(flota[j].getNombre(), actual.getNombre())){
        flota[j + 1] = flota[j];
        j--;
      }
      flota[j + 1] = actual;
    }
  }
\end{lstlisting}
\begin{itemize}
    \item Estamos usando una variable para almacenar nuestro elemento a ubicar
    \item Luego recorremos tantos elementos a la derecha como elementos menores a nuestro pivote hayan
    \item Y finalmente colocamos nuestro pivote en la posicion que queda libre
\end{itemize}
\subsubsection{Ejecucion}
\begin{itemize}
    \item Ahora veremos la prueba del metodo:
\end{itemize}
    \begin{lstlisting}[language=bash, caption={Prueba del metodo}]
    Ahora usaremos el algoritmo de ordenamiento seleccion, con respecto a Nombres.
    Naves ordenadas:
    
    
    Nave 1:
    Nombre: Alvaro
    Estado: true
    Puntos: 6
    
    Nave 2:
    Nombre: Bannia
    Estado: true
    Puntos: 89
    
    Nave 3:
    Nombre: cesar
    Estado: true
    Puntos: 23
    
      
\end{lstlisting}





\subsection{Ordenamiento por puntos: insercion}
\begin{itemize}
    \item Ahora usaremos este mismo algoritmo para ordenar el array con respecto a los puntos
    \item Es la misma estructura y el mismo sentido, en lugar de unsar getNombre, usamos getPuntos
\end{itemize}
\begin{lstlisting}[language=java, caption={Metodo para ordenar por puntos}]
      public static void ordenarPorPuntosInsercion(Nave[] flota){
    for(int i = 1; i < flota.length; i++){
      Nave actual = flota[i];
      int j = i - 1;
      while(j >= 0 && flota[j].getPuntos() > actual.getPuntos()){
        flota[j + 1] = flota[j];
        j--;
      }
      flota[j + 1] = actual;
    }
  }
\end{lstlisting}
\subsubsection{Ejecucion}
\begin{itemize}
    \item Ahora veremos la prueba del metodo:
\end{itemize}
    \begin{lstlisting}[language=bash, caption={Prueba del metodo}]
    Ahora usaremos el algoritmo de ordenamiento insercion, con respecto a Puntos.
    Naves ordenadas:
    
    
    Nave 1:
    Nombre: Alvaro
    Estado: true
    Puntos: 6
    
    Nave 2:
    Nombre: cesar
    Estado: true
    Puntos: 23
    
    Nave 3:
    Nombre: Bannia
    Estado: true
    Puntos: 89      
\end{lstlisting}



\subsection{Busqueda binaria}
\begin{itemize}
    \item La busqueda binaria es una forma mas eficiente de buscar un elemento en una lista
    \item Como pre requisito, es que nuestro arreglo tiene que estar ordenado alfabeticamente, para eso usamos cualquier algoritmo de ordenamiento por nombre visto anteriormente
    \item El tiempo de complejidad es o(log n), lo que es muy bueno
\end{itemize}
\begin{lstlisting}[language=java, caption={Metodo para hacer una busqueda binaria}]
      public static int busquedaBinariaNombre(Nave[] flota, String s){
    ordenarPorNombreSeleccion(flota);
    int min = 0;
    int max = flota.length - 1;
    while(min <= max){
      int medio = (max + min) / 2;
      if(flota[medio].getNombre().equalsIgnoreCase(s))
        return medio;
      if(esMayor(s, flota[medio].getNombre()))
        min = medio + 1;
      else
        max = medio - 1;
    }
    return - 1;
  }
\end{lstlisting}
\subsubsection{Ejecucion}
\begin{itemize}
    \item Ahora veremos la prueba del metodo:
\end{itemize}
    \begin{lstlisting}[language=bash, caption={Prueba del metodo}]
    Busqueda binaria, inserte un nombre:
    cesar
    Nombre: cesar
    Estado: true
    Puntos: 23
\end{lstlisting}



\section{Commits importantes}
\begin{itemize}
    \item Ahora veremos el registro de los commits principales:
\end{itemize}
\begin{lstlisting}

commit 9b1f5460554796ac356fd9dcae4fc20bb3d3a123
Author: JhonatanDczel <jariasq@unsa.edu.pe>
Date:   Wed Sep 27 01:46:26 2023 -0500

    Actividad 1: Implementacion del metodo de ordenamiento de Puntos por insercion

commit a68dcb2c45cf068ed3fdf3733bfd3113f16abe6c
Author: JhonatanDczel <jariasq@unsa.edu.pe>
Date:   Wed Sep 27 01:18:23 2023 -0500

    Actividad 1: Implementacion del algoritmo de ordenamiento de puntos por seleccion

commit 18c3176724a4529cf0e77038215cdf5e180088ae
Author: JhonatanDczel <jariasq@unsa.edu.pe>
Date:   Wed Sep 27 01:07:19 2023 -0500

    Actividad 1: Implementacion del metodo de busqueda binaria por nombre

commit 0db1011f08fff2f98effb4123c5d2136ea68879d
Author: JhonatanDczel <jariasq@unsa.edu.pe>
Date:   Tue Sep 26 23:54:17 2023 -0500

    Actividad 1: Implementacion del metodo de ordenamiento por puntos usando el algoritmo burbuja

\end{lstlisting}
\begin{itemize}
    \item Cada commit representa el uso de un algoritmo diferente
    \item Esta el algoritmo burbuja, insercion y seleccion, asi como la busqueda binaria
\end{itemize}
%Rubricas de evaluacion
	\clearpage
 	\section{\textcolor{red}{Rúbricas}}
	\subsection{\textcolor{red}{Rúbrica para el contenido del Informe y demostración}}
	\begin{itemize}			
		\item El alumno debe marcar o dejar en blanco en celdas de la columna \textbf{Checklist} si cumplio con el ítem correspondiente.
		\item El alumno debe autocalificarse en la columna \textbf{Estudiante} de acuerdo a la siguiente tabla:
	
		\begin{table}[ht]
			\caption{Niveles de desempeño}
			\begin{center}
			\begin{tabular}{ccccc}
    			\hline
    			 & \multicolumn{4}{c}{Nivel}\\
    			\cline{1-5}
    			\textbf{Puntos} & Insatisfactorio 25\%& En Proceso 50\% & Satisfactorio 75\% & Sobresaliente 100\%\\
    			\textbf{2.0}&0.5&1.0&1.5&2.0\\
    			\textbf{4.0}&1.0&2.0&3.0&4.0\\
    		\hline
			\end{tabular}
		\end{center}
	\end{table}	
	
	\end{itemize}
	
	\begin{table}[H]
		\caption{Rúbrica para contenido del Informe y demostración}
		\setlength{\tabcolsep}{0.5em} % for the horizontal padding
		{\renewcommand{\arraystretch}{1.5}% for the vertical padding
		%\begin{center}
		\begin{tabular}{|p{2.7cm}|p{7cm}|x{1.3cm}|p{1.2cm}|p{1.5cm}|p{1.1cm}|}
			\hline
    		\multicolumn{2}{|c|}{Contenido y demostración} & Puntos & Checklist & Estudiante & Profesor\\
			\hline
			\textbf{1. GitHub} & Hay enlace URL activo del directorio para el  laboratorio hacia su repositorio GitHub con código fuente terminado y fácil de revisar. &2 &X &2 & \\ 
			\hline
			\textbf{2. Commits} &  Hay capturas de pantalla de los commits más importantes con sus explicaciones detalladas. (El profesor puede preguntar para refrendar calificación). &4 &X &4 & \\ 
			\hline 
			\textbf{3. Código fuente} &  Hay porciones de código fuente importantes con numeración y explicaciones detalladas de sus funciones. &2 &X &2 & \\ 
			\hline 
			\textbf{4. Ejecución} & Se incluyen ejecuciones/pruebas del código fuente  explicadas gradualmente. &2 &X &1.5 & \\ 
			\hline			
			\textbf{5. Pregunta} & Se responde con completitud a la pregunta formulada en la tarea.  (El profesor puede preguntar para refrendar calificación).  &2 &X &2 & \\ 
			\hline	
			\textbf{6. Fechas} & Las fechas de modificación del código fuente estan dentro de los plazos de fecha de entrega establecidos. &2 &X &2 & \\ 
			\hline 
			\textbf{7. Ortografía} & El documento no muestra errores ortográficos. &2 &X &1.5 & \\ 
			\hline 
			\textbf{8. Madurez} & El Informe muestra de manera general una evolución de la madurez del código fuente,  explicaciones puntuales pero precisas y un acabado impecable.   (El profesor puede preguntar para refrendar calificación).  &4 &X &4 & \\ 
			\hline
			\multicolumn{2}{|c|}{\textbf{Total}} &20 & &19 & \\ 
			\hline
		\end{tabular}
		%\end{center}
		%\label{tab:multicol}
		}
	\end{table}
	

\clearpage
	
%\clearpage
%\bibliographystyle{apalike}
%\bibliographystyle{IEEEtranN}
%\bibliography{bibliography}
			
\end{document}
